%TC:ignore
% ============================================================
%  packages.tex – Pakete und Layout-Konfiguration
%  Orientiert am FOM Leitfaden IT/ING 2024
% ============================================================

% ── Sprache, Encoding, Schrift ────────────────────────────────
\usepackage[utf8]{inputenc}
\usepackage[T1]{fontenc}
\usepackage[ngerman]{babel}
\usepackage{csquotes}            % korrekte Anf\"uhrungszeichen (deutsch: „…")

% Deutsche Bezeichnungen f\"ur \autoref
\addto\extrasngerman{%
  \renewcommand{\chapterautorefname}{Kapitel}%
  \renewcommand{\sectionautorefname}{Abschnitt}%
  \renewcommand{\subsectionautorefname}{Abschnitt}%
  \renewcommand{\subsubsectionautorefname}{Abschnitt}%
  \renewcommand{\figureautorefname}{Abbildung}%
  \renewcommand{\tableautorefname}{Tabelle}%
}

% ── Schriftart: Arial-Ersatz (Helvetica) in 11pt ─────────────
\usepackage[scaled]{helvet}
\renewcommand{\familydefault}{\sfdefault}

% ── Seitenr\"ander nach Leitfaden \S5.2.1 ────────────────────
\usepackage[
  a4paper,
  top=4cm,
  bottom=2cm,
  left=4cm,
  right=2cm,
  headsep=1.4cm,
  headheight=17pt,
  footskip=1.5cm,
]{geometry}

% ── Zeilenabstand / Absatz ────────────────────────────────────
\usepackage{setspace}
\onehalfspacing                  % 1,5-zeilig
\setlength{\parskip}{6pt}
\setlength{\parindent}{0pt}

\usepackage{microtype}

% ── Grafiken ──────────────────────────────────────────────────
\usepackage{graphicx}
% Floats duerfen nicht ueber Abschnittsgrenzen hinweg abdriften
\usepackage[section]{placeins}

% Grosszuegigere Float-Platzierung: Bilder/Tabellen bleiben im Fliesstext
% statt auf eine eigene Float-Seite verbannt zu werden
\renewcommand{\topfraction}{0.85}
\renewcommand{\bottomfraction}{0.7}
\renewcommand{\textfraction}{0.15}
\renewcommand{\floatpagefraction}{0.66}
\setcounter{topnumber}{3}
\setcounter{bottomnumber}{2}
\setcounter{totalnumber}{4}

% Float-Seiten: Inhalt oben ausrichten statt vertikal zu zentrieren
\makeatletter
\setlength{\@fptop}{0pt}
\setlength{\@fpbot}{0pt plus 1fil}
\makeatother

% Hurenkinder-/Schusterjungen-Schutz: keine einzelnen Zeilen eines
% Absatzes allein am Seitenanfang oder -ende
\clubpenalty=10000
\widowpenalty=10000
\displaywidowpenalty=10000
% ── URLs ──────────────────────────────────────────────────────
\usepackage{xurl}
\urlstyle{same}

% ── Tabellen ──────────────────────────────────────────────────
\usepackage{longtable}
\usepackage{booktabs}
\usepackage{array}
\usepackage{ragged2e}
\usepackage{tabularx}
\usepackage{eurosym}

\newcolumntype{Y}{>{\RaggedRight\arraybackslash}X}

% ── Mathematik ────────────────────────────────────────────────
\usepackage{amsmath}
\usepackage{amssymb}

% ── Fortlaufende Nummerierung: 1, 2, 3 statt 2.1, 2.2 ────────
\usepackage{chngcntr}
\counterwithout{figure}{chapter}
\renewcommand{\thefigure}{\arabic{figure}}
\counterwithout{table}{chapter}
\renewcommand{\thetable}{\arabic{table}}

% ── "Abbildung 1 …" / "Tabelle 1 …" im Verzeichnis ──────────
\newcommand*{\tocfigureformat}[1]{Abbildung~#1}
\newcommand*{\toctableformat}[1]{Tabelle~#1}

\DeclareTOCStyleEntry[
  indent=0pt,
  numwidth=2.4cm,
  entrynumberformat=\tocfigureformat,
]{tocline}{figure}

\DeclareTOCStyleEntry[
  indent=0pt,
  numwidth=2.4cm,
  entrynumberformat=\toctableformat,
]{tocline}{table}

% ── Kopf-/Fu\3zeile nach Leitfaden ───────────────────────────
\usepackage[automark]{scrlayer-scrpage}
\clearpairofpagestyles
\automark[section]{chapter}
\renewcommand*\chaptermarkformat{}
\ihead{\leftmark}
\ohead{\pagemark}
\setheadsepline{0.4pt}
\renewcommand*\chapterpagestyle{scrheadings}
\setkomafont{pagehead}{\normalfont\normalsize}
\setkomafont{pagination}{\normalfont\normalsize}

% ── Kapitelstart ohne Extra-Abstand ──────────────────────────
\makeatletter
\renewcommand*{\chapterheadstartvskip}{\vspace*{0pt}}
\makeatother

% ── \"Uberschriften-Abst\"ande nach Leitfaden ────────────────
\RedeclareSectionCommand[beforeskip=12pt, afterskip=6pt]{chapter}
\RedeclareSectionCommand[beforeskip=12pt, afterskip=6pt]{section}
\RedeclareSectionCommand[beforeskip=12pt, afterskip=6pt]{subsection}
\RedeclareSectionCommand[beforeskip=12pt, afterskip=6pt]{subsubsection}

% ── Abbildungen/Tabellen: Beschriftung ───────────────────────
\usepackage[
  font=small,
  labelfont=bf,
  format=plain,
  justification=centering,
  aboveskip=12pt,
  belowskip=6pt
]{caption}

% ── Glossaries-extra (Abk\"urzungen) ─────────────────────────
\usepackage[
  acronym,
  automake,
  nopostdot,
  nonumberlist,
  toc,
  section=chapter,
  nogroupskip,
]{glossaries-extra}
\makeglossaries

\setabbreviationstyle[acronym]{long-short}

\newglossarystyle{abbrv-fom}{%
  \renewenvironment{theglossary}%
    {\begin{longtable}[l]{@{}p{2cm}@{\hspace{0.2cm}}p{\dimexpr\linewidth-2cm\relax}@{}}}%
    {\end{longtable}}%
  \renewcommand*{\glossaryheader}{}%
  \renewcommand*{\glsgroupheading}[1]{}%
  \renewcommand*{\glsgroupskip}{}%
  \renewcommand*{\glossentry}[2]{%
    \glstarget{##1}{\glossentryname{##1}} &
    \glossentrydesc{##1}\tabularnewline
  }%
  \renewcommand*{\subglossentry}[3]{}%
}

% ============================================================
%  LITERATURVERZEICHNIS – BibLaTeX (APA), FOM-Leitfaden
% ============================================================
\usepackage[
  backend=biber,
  style=apa,
  autocite=inline,
  maxcitenames=2,
  maxbibnames=20,
  giveninits=false,
  uniquename=false,
  uniquelist=false,
  dashed=false,
  sorting=anyt,
  urldate=long,
  dateabbrev=false,
]{biblatex}

\DeclareFieldFormat{pages}{S.~#1}
\addbibresource{literature.bib}

\DefineBibliographyStrings{ngerman}{%
  nodate    = {o\adddot\space D\adddot},
  retrieved = {Abgerufen\space am},
  from      = {von},
}
%TC:endignore
